\documentclass[10pt,a4paper]{amsart}
\usepackage[margin=19mm]{geometry}
\usepackage{amsmath,amssymb,mathtools}
\usepackage[scr=rsfs,cal=euler]{mathalfa}
\usepackage{xfrac}
\usepackage[unicode,hidelinks,bookmarks=false]{hyperref}
\newcommand{\ie}{i.e.\ }
\newcommand{\cf}{cf.\ }
\newcommand{\resp}{respectively\ }
\newcommand{\Subsection}{\subsection}
\providecommand{\nobreakdash}{\nobreak\hbox{-}}
\setlength{\emergencystretch}{2em}
\multlinegap0pt
\usepackage{bm}
\usepackage{bbm}
\usepackage{dsfont}
\usepackage{xspace}
\usepackage{cancel}
\usepackage{mathtools}
\usepackage{amsthm}
\makeatletter
\@ifundefined{theorem}{\newtheorem{theorem}{Theorem}}{}
\@ifundefined{proposition}{\newtheorem{proposition}[theorem]{Proposition}}{}
\@ifundefined{remark}{\newtheorem{remark}[theorem]{Remark}}{}
\makeatother
\renewcommand{\geq}{\geqslant}
\renewcommand{\leq}{\leqslant}
\title[Released packing functions in graphs]{Released packing functions in graphs}
\author{Pablo Fekete, Erica Hinrichsen, Valeria Leoni,, Mar\'ia In\'es Lopez Pujato}
\date{Source: arXiv:2608.11169v1 (CC BY 4.0)}
\begin{document}
\maketitle

\noindent\emph{Source and licence.} Released packing functions in graphs. Version arXiv:2608.11169v1 (CC BY 4.0), distributed under the Creative Commons Attribution 4.0 International licence, as recorded in the arXiv licence metadata for this exact version and shown on the abstract page https://arxiv.org/abs/2608.11169v1. Full rights notice: licenses/arxiv-2608.11169-CC-BY-4.0.txt.\par\smallskip

\noindent\textit{Editorial scope.} Counted unit: the Theorem whose source label is \texttt{th:4.7} in the article (source0.tex lines 469--477, source label \texttt{th:4.7}), delivered together with the article's own complete original proof of it (source0.tex lines 479--501, one \texttt{proof} environment of 23 lines). No mathematics is written, repaired or re-derived: statement and proof are the article's own text line for line. Required context retained uncounted: prop:dimensionofP (lines 394--397); eqn: vecindad con u=1 (lines 418--420); rem:V1 (lines 427--429). Every definition, display and label of those lines is retained unchanged. Omitted from this package and not redistributed: 1--393, 398--417, 421--426, 430--468, 478--478, 502--547. Nothing inside a retained excerpt is omitted or altered: every retained body is delivered exactly as the article's lines are written, so no omission receipt of any kind accompanies this package. Citations: the retained excerpts carry no citation command at all, so no bibliography entry of the article is transcribed and no reference is redistributed. Reference adaptation: none was needed and none was made; every reference inside the delivered unit resolves inside it, so no unresolved reference or citation is produced. Printed slips kept literal: no printed word, symbol, hypothesis or number of the counted statement or of its proof was altered. Layout and class: the article's own class is \texttt{article} with packages including inputenc, indentfirst, csquotes, keytheorems, zref-clever, soul, geometry, layout, amsmath, amssymb, amsthm, amsfonts, comment, graphicx; the wrapper is the lane's pinned \texttt{amsart} editorial wrapper at 10pt on A4, which restates the theorem-like environments the retained text needs and adds the article's own macro definitions that the retained text needs (\texttt{\textbackslash geq}, \texttt{\textbackslash leq}), with extra wrapper packages (bm, bbm, dsfont, xspace, cancel, mathtools). The delivered numbering is the unit's own. Assets: no retained span contains a figure environment, a graphics-inclusion command, a PDF-page inclusion, a table or a TikZ picture; no image file is copied into this package, the package declares no asset dependency and no image file is redistributed with it. Licence: version arXiv:2608.11169v1 of this article is distributed under the Creative Commons Attribution 4.0 International licence, as recorded in the arXiv licence metadata for this exact version and shown on the abstract page https://arxiv.org/abs/2608.11169v1. A full rights notice is delivered at licenses/arxiv-2608.11169-CC-BY-4.0.txt. Characters: every retained excerpt is pure ASCII, so the delivered engine, XeLaTeX with its default Latin Modern text font, reports no missing character.
\begin{proposition}\label{prop:dimensionofP}
Given a graph $G=(V,E)$ and vectors $\mathbf{u}\geq \mathbf{0}$ and $\mathbf{k}\geq \mathbf{1}$ such that  $u_v\leq k_v\leq u(N_G[v])$  for all $v\in V$, let  $V_0(\mathbf{u})=\{v\in V:\ u_v=0\}$. Then, the dimension of $P(G,\mathbf{k},\mathbf{u})$, 
$dim(P(G,\mathbf{k},\mathbf{u}))$, is equal to $|V|-|V_0(\mathbf{u)}|$.
\end{proposition}

\begin{align}
(d_G(v)+1 - k_v)x_v + x( N_G(v))  \leq d_G(v) &\quad (v\in V\text{ with } k_v\leq d_G(v) ) \label{eqn: vecindad con u=1} \\ 
0\leq x_v \leq 1 \label{eqn:0,1} &\quad (v\in V).\end{align}

\begin{remark}\label{rem:V1}
   For a graph $G=(V,E)$, $\mathbf{u}=\mathbf{1}$, $\mathbf{k}\geq \mathbf{1}$ and  $\mathbf{x}\in \{0,1\}^{|V|}$, let $$V_1(\mathbf{x})=\{v\in V:\ x_v=1\}.$$ Then $\mathbf{x}\in P(G,\mathbf{k},\mathbf{1})$ if and only if $|V_1(\mathbf{x})\cap N_G(v)|\leq k_v-1$ for all $v\in V_1(\mathbf{x})$.
\end{remark}

\begin{theorem}\label{th:4.7}
    Given a graph $G=(V,E)$ and  $v\in V$ with $d_G(v)\geq 2$, inequality in \eqref{eqn: vecindad con u=1} for $v$ defines a facet of $P(G,2\cdot\mathbf{1},\mathbf{1})$ if and only if the following conditions hold:
    \begin{enumerate}
        \item[i.] $\Delta(G[N_G(v)])\leq 1$, and
        \item[ii.] $N_G(v)\setminus N_G(v')\neq \emptyset$, for all $v'\in V\setminus N_G[v]$.
    \end{enumerate}
    
     
\end{theorem}

\begin{proof}
    To simplify notation, we denote $P=P(G,2\cdot\mathbf{1},\mathbf{1})$. Let $\mathcal{F}$ be the face of $P$ defined by  inequality in \eqref{eqn: vecindad con u=1}, i.e., $\mathcal{F}=P\cap H$, where:
    \[H=\{\mathbf{x}:\ (d_G(v)-1)x_v +x(N_G(v)) = d_G(v) \}.\]

    
    Suppose $\mathcal{F}$ is a facet of $P$. As $P$ is a full dimensional polytope by Proposition \ref{prop:dimensionofP}, $dim(\mathcal{F})=|V|-1$. This implies that $\mathcal{F}$ is not contained in any hyperplane $H_{w,b}:=\{\mathbf{x}: x_w=b\}$, for $w\in V$ and $b\in\{0,1\}$, since otherwise $dim(\mathcal{F})\leq |V|-2$.
    
    We first prove the necessity of condition \textit{i.} By contradiction, assume $\Delta(G[N(v)])\geq 2$. Then there exist  $w_1,w_2,w_3\in N(v)$ satisfying 
    $w_1w_2,w_1w_3\in E$. Let $\mathbf{x}\in H\cap \{0,1\}^{|V|}$ such that $x_v=0$. Then $x_w=1$ for all $w\in N_G(v)$. In particular, $w_1\in V_1(\mathbf{x})$ and $\{w_2,w_3\}\subseteq V_1(\mathbf{x})\cap N_G(w_1)$. By Remark~\ref{rem:V1}, $\mathbf{x}\notin P$. Then, $x_v=1$ for all  $\mathbf{x}\in \mathcal{F}$ and thus $\mathcal{F}\subseteq H_{v,1}$, contradicting the fact that $\mathcal{F}$ is a facet. Hence, $\Delta(G[N_G(v)])\leq 1$.

    To prove the necessity of condition \textit{ii.} suppose, by contradiction, that there is a vertex $v'\in V\setminus N_G[v]$ such that $N_G(v)\subseteq N_G(v')$ and let $\mathbf{x}\in H\cap \{0,1\}^{|V|}$ with $x_{v'}=1$. If $x_v=1$, then there is (exactly) one vertex $w\in N_G(v)$ such that $x_w=1$, and then we have $\{v,v'\}\subseteq V_1(\mathbf{x})\cap N_G(w)$. Else ($x_v=0$) we have $x_w=1$ for all $w\in N_G(v)$. This implies $|V_1(\mathbf{x})\cap N_G(v')|\geq 2$, as $d_G(v)\geq 2$ and $N_G(v)\subseteq N_G(v')$. In either case, Remark~\ref{rem:V1} implies $\mathbf{x}\notin P$. Hence, $\mathcal{F}\subseteq H_{v', 0}$, contradicting the fact that $\mathcal{F}$ is a facet of $P$. Therefore, $N_G(v)\setminus N_G(v')\neq \emptyset$, for all $v'\in V\setminus N_G[v]$.

    Finally, we prove that $\mathcal{F}$ is a facet if conditions \textit{i.} and \textit{ii.} hold. 
    On the one hand, \textit{ii.} implies that for each $v'\in V\setminus N_G[v]$ there exists  $w'\in N_G(v)\setminus N_G(v')$. Consider the points $\mathbf{x}^{v'}\in\{0,1\}^{|V|}$ given by $\mathbf{x}^{v'}=\mathbf{e}^v+\mathbf{e}^{w'}+\mathbf{e}^{v'}$, for $v'\in V\setminus N_G[v]$, together with the points $\mathbf{x}^v=\sum_{w\in N_G(v)}\mathbf{e}^w$ and $\mathbf{x}^{w}=\mathbf{e}^v+\mathbf{e}^w$, for each $w\in N_G(v)$. To prove  that these $|V|$ points are affinely independent points we proceed as follows. Suppose $\sum_i\lambda_i \mathbf{x}^{i}= \mathbf{0}$ with
$\sum_i\lambda_i=0$, where the points are $\mathbf{x}^{v}$, $\{\mathbf{x}^{w}\}_{w\in N(v)}$ and
$\{\mathbf{x}^{v'}\}_{v'\in V\setminus N[v]}$. For each $v'\in V\setminus N_G[v]$, component
$v'$ equals $1$ only in $\mathbf{x}^{v'}$, then $\lambda_{v'}=0$. Besides, component $v$ is $0$ in
$\mathbf{x}^{v}$ and $1$ in every $\mathbf{x}^{w}$, hence $\sum_{w\in N_G(v)}\lambda_{w}=0$, and then
$\lambda_v=0$. Finally, for each $w\in N_G(v)$, component $w$ is $1$ only in
$\mathbf{x}^{w}$ among the remaining points, so $\lambda_w=0$. Thus all coefficients
vanish and the points exhibited are affinely independent. Clearly these points belong to $H$, and from Remark~\ref{rem:V1} and \textit{i.}, they all belong to $P$. Hence, $\mathcal{F}$ is a facet of $P$.

\end{proof}

\end{document}
